\documentclass[11pt]{article}

%==============================================================================
%  The Local-Chart Simplicial Complex as an Organizing Substrate for LPS
%  A cross-cutting design note: the nerve of the chart cover as the native
%  combinatorial substrate for conditional-expectation and local-covariance
%  estimation. Organizational (not topological) view. All figures original.
%==============================================================================

\usepackage[margin=1in]{geometry}
\usepackage{amsmath,amssymb,amsthm,mathtools}
\usepackage{graphicx}
\usepackage{xcolor}
\usepackage{booktabs}
\usepackage{enumitem}
\usepackage{microtype}
\usepackage{caption}
\usepackage[most]{tcolorbox}
\usepackage{datetime2}
\DTMsetup{showzone=false}     % timestamp in US Eastern time (build sets TZ); see latexmkrc
\usepackage[hidelinks]{hyperref}
\usepackage[section]{placeins}
\widowpenalty=10000
\clubpenalty=10000
\raggedbottom

\definecolor{ink}{HTML}{24303D}
\definecolor{accent}{HTML}{2563EB}
\definecolor{accent2}{HTML}{B45309}
\captionsetup{font=small,labelfont=bf}
\setlist{topsep=3pt,itemsep=2pt,parsep=1pt}
\theoremstyle{definition}
\newtheorem{definition}{Definition}[section]

\newtcolorbox{plain}{
  enhanced, colback=accent!4, colframe=accent!55!black, boxrule=0.4pt, arc=2pt,
  left=7pt, right=7pt, top=4pt, bottom=4pt, title=In plain terms,
  fonttitle=\bfseries\sffamily\small, coltitle=accent!50!black,
  attach title to upper={\quad}}
\newtcolorbox{caution}{
  enhanced, colback=accent2!5, colframe=accent2!60!black, boxrule=0.4pt, arc=2pt,
  left=7pt, right=7pt, top=4pt, bottom=4pt, title=A caution worth stating,
  fonttitle=\bfseries\sffamily\small, coltitle=accent2!55!black,
  attach title to upper={\quad}}

\newcommand{\R}{\mathbb{R}}
\newcommand{\E}{\mathbb{E}}
\newcommand{\nerve}{\mathcal{N}}
\newcommand{\fig}[3]{%
  \begin{figure}[htbp]\centering
  \includegraphics[width=#2]{figures_ldr/#1}%
  \caption{#3}\label{fig:#1}\end{figure}}

\title{\textbf{The Local-Chart Simplicial Complex as an Organizing Substrate
for LPS}\\[4pt]
\large From the adaptive kNN graph to the nerve of the chart cover --- an
organizational (not topological) scaffold for refining conditional-expectation
and local-covariance estimation}
\author{\small Internal design note --- LPS / \texttt{geosmooth} program\\[2pt]
\footnotesize source: \nolinkurl{/Users/pgajer/current_projects/geosmooth/project_briefs/}\\[1pt]
\footnotesize\nolinkurl{lps_chart_nerve_substrate_2026-06-10.tex}}
\date{\DTMnow\ \textnormal{\footnotesize ET}}

\begin{document}
\maketitle

\begin{abstract}\noindent
Every estimator in the LPS program currently sits on top of a graph that is built
\emph{before, and independently of}, the estimation problem: an adaptive-radius
$k$-nearest-neighbor graph. This note argues that once we commit to \emph{local
chart} models, the natural combinatorial substrate is not that input graph but the
\emph{nerve of the chart cover} --- the simplicial complex whose vertices are the
charts and whose simplices record which charts overlap. The thesis is
deliberately organizational, not topological: we are not interested in
reconstructing the data's homotopy type, only in a faithful bookkeeping of
\emph{which local estimates should agree, and how strongly}, so that conditional
expectations and local covariances can be refined by gluing. Seen this way, the
program already uses the shadow of this object --- PS-LPS's overlap coupling and
the chart-aware association layer's overlap weights are exactly the nerve's
$1$-skeleton. Making the nerve native buys three things a kNN graph cannot give:
variable chart sizes (an adaptive cover), higher-order overlaps (constraints where
three or more charts meet), and a substrate whose quality is judged by held-out
estimation accuracy rather than fixed by an a-priori rule. We lay out what each
LPS layer gains, the one criterion that makes a complex ``good'' for our purposes
(a statistical one), the compositional case where part of the complex is known for
free, the real costs, and a minimal-first path keyed to the synthetic generators
of the implementation notes.
\par\medskip\noindent\footnotesize\textit{On the figures.} All figures are
original schematics drawn for this note.
\end{abstract}

{\small\tableofcontents}

%==============================================================================
\section{The substrate, and a deliberate restriction of ambition}\label{sec:thesis}
%==============================================================================

The LPS pipeline builds, in order: a neighbor graph on the samples; a local chart
at each anchor (local PCA, with an estimated local dimension); a local polynomial
fit on that chart; and, in PS-LPS and the chart-aware association layer, a
coupling that ties \emph{overlapping} charts together. The first object --- the
adaptive-radius kNN graph --- is built by a generic proximity rule that knows
nothing about the regression surface or the covariance structure we are trying to
estimate. It is scaffolding.

The thesis of this note is that the charts themselves induce a better substrate.
A collection of charts is a \emph{cover} of the data: each chart's support is a
subset $U_c$ of samples, and together the $U_c$ blanket the cloud. The
combinatorial object that records how a cover fits together is its \emph{nerve}.

\begin{definition}[Nerve of the chart cover]\label{def:nerve}
Let $\{U_c\}_{c}$ be the supports of the local charts. The nerve $\nerve$ is the
simplicial complex with one vertex per chart $c$, an edge $\{c,c'\}$ whenever
$U_c\cap U_{c'}\neq\varnothing$, a triangle $\{c,c',c''\}$ whenever
$U_c\cap U_{c'}\cap U_{c''}\neq\varnothing$, and so on. Each simplex carries the
size of the corresponding overlap as a weight.
\end{definition}

\fig{fig_substrate.pdf}{\linewidth}{The shift in substrate. (a) The input
scaffold: an adaptive-radius kNN graph, pairwise and fixed before estimation.
(b) The charts form a \emph{cover} whose pieces can vary in size with the local
geometry. (c) The nerve of that cover: charts are nodes, an edge marks a pairwise
overlap, and a filled triangle marks a place where three charts overlap at once
--- structure a pairwise graph cannot record.}

We must be precise about what we want from $\nerve$, because the word ``nerve''
comes loaded with topology. The classical reason to care about a nerve is the
Nerve Theorem: for a sufficiently nice (``good'') cover, the nerve is homotopy
equivalent to the union of the cover, so one can read the space's topology off a
finite complex. \emph{We are explicitly not relying on that.} The goal here is
biology, and the use is organizational: $\nerve$ is a faithful ledger of which
local estimates overlap, and therefore of which ones should be made to agree and
with what strength. Whether $\nerve$ reproduces the data's homology is beside the
point; what matters is whether gluing local estimates over $\nerve$ produces
better conditional expectations and covariances. That reframing --- from a
topological criterion to a statistical one --- runs through the whole note.

\begin{plain}
Think of the charts as overlapping photographs that together cover a landscape.
The ``nerve'' is just the index card that says which photos overlap, and which
three overlap at a corner. We are not trying to rebuild the landscape's shape from
the index card; we are using the card to know which photos must be stitched to
agree where they overlap --- so the stitched panorama (our estimate of the
surface) is seamless.
\end{plain}

%==============================================================================
\section{The program already uses the nerve's shadow}\label{sec:shadow}
%==============================================================================

This is not an exotic import. Two pieces of the existing program are, up to
naming, already defined on $\nerve$.

\begin{itemize}[leftmargin=1.4em]
\item \textbf{PS-LPS} couples two charts when their supports share points and
penalizes disagreement of their predictions on the shared points. ``Charts that
share support points'' is precisely an \emph{edge of $\nerve$}; the shared points
are the overlap that weights it. PS-LPS is a synchronization defined on the nerve's
$1$-skeleton.
\item \textbf{The chart-aware association layer} fuses neighboring charts'
covariance/precision blocks with an overlap weight $\zeta_{cc'}=\sum_j\min(w_{cj},
w_{c'j})$ --- again an edge weight of $\nerve$, measuring how much two charts
overlap.
\end{itemize}

So the program is already doing nerve-based gluing; it just derives the adjacency
from the kNN graph and stops at pairwise edges. The proposal is to (i) derive
$\nerve$ \emph{natively} from the charts (so it adapts as the charts adapt), and
(ii) use more of it than the bare $1$-skeleton.

%==============================================================================
\section{Why the nerve organizes better than any kNN graph}\label{sec:why}
%==============================================================================

Three differences matter, and all three are in service of estimation, not topology.

\paragraph{Variable chart size $=$ an adaptive cover.} A kNN graph fixes the arity
(or radius) of every neighborhood before looking at the problem. Charts need not:
a chart can be large where the data is sparse or locally one-dimensional and small
where it is dense or higher-dimensional. The nerve of a variable-size cover encodes
this adaptivity directly --- two regions are coupled exactly when their
\emph{charts} (sized to the local geometry) overlap, not when an arbitrary fixed
$k$ says so.

\paragraph{Higher-order overlaps $=$ constraints a graph cannot express.} Where
three charts overlap, the right object is a single three-way agreement, not three
separate pairwise ones. A pairwise graph can only ever say ``$a$ agrees with $b$,
$b$ with $c$, $a$ with $c$''; it cannot say ``$a$, $b$, $c$ agree on their common
region.'' The latter is a genuinely different (and at junctions, more appropriate)
constraint, and it is exactly a $2$-simplex of $\nerve$ (Figure~\ref{fig:fig_glue.pdf},
right). This is the same reason the clique/Rips complex induced by a kNN graph ---
which manufactures higher simplices from pairwise edges --- is a poorer record of
the cover than the nerve, which records the actual multi-way intersections.

\paragraph{The substrate becomes part of the model.} Because $\nerve$ is built
from the charts, and the charts are chosen to fit the estimation problem, the
substrate is no longer fixed input: it is selected, and can be judged by how well
estimates glued over it predict held-out data (Section~\ref{sec:criterion}).

\begin{caution}
None of this requires the cover to be ``good'' in the topological sense. A kNN
graph and the nerve will sometimes disagree about whether two regions should be
coupled, and near a stratum boundary the charts may overlap in ways that would
\emph{not} make a topologically faithful nerve. For our purposes that is fine:
an extra or missing edge changes which estimates are tied together and how
strongly --- a modeling choice to be validated by accuracy, not a topological
error. We keep the nerve as a statistical ledger and let held-out error referee it.
\end{caution}

%==============================================================================
\section{What each LPS layer gains}\label{sec:gains}
%==============================================================================

\subsection{Conditional expectation (PS-LPS)}
Moving PS-LPS onto the native nerve changes three things. \emph{Adaptivity:} the
coupling graph is the nerve of variable-size charts, so two anchors are
synchronized exactly when their (geometry-sized) charts overlap, with strength set
by the overlap. \emph{Higher order:} on a $2$-simplex, synchronize the three
charts' predictions on their common overlap as a single constraint --- the natural
thing at a junction where three local models meet
(Figure~\ref{fig:fig_glue.pdf}). \emph{Honest weights:} edge and triangle weights
come from actual intersection sizes, not from a kNN kernel evaluated at fixed
arity.

\fig{fig_glue.pdf}{\linewidth}{Gluing on overlaps. Left: where two charts overlap
(an edge of $\nerve$) their local fits --- both the predicted surface and the local
covariance --- should agree on the shared region. Right: where three charts overlap
(a $2$-simplex) there is a three-way agreement that a pairwise graph cannot
express. These overlaps are exactly where, and how strongly, local estimates are
tied together.}

\subsection{Local covariance and association}
The chart-aware layer estimates a covariance/precision per chart and fuses
neighbors. The overlap of two charts is precisely the set of samples on which
their two covariance estimates are looking at the same data --- so $\nerve$ tells
us \emph{where} the covariance estimates should agree and \emph{how strongly}
(the overlap weight). Variable chart size lets the covariance model widen its
support where data is scarce. And the higher simplices give three-way covariance
agreement at junctions, where the support layer (which taxa co-occur) reorganizes.

\subsection{The local-dimension / chart-size field}
The dimension field of the implementation notes minimizes
$\sum_i c_i(d_i)+\lambda\sum_{(i,j)}w_{ij}\rho(d_i,d_j)$ over a graph; that graph
should be the nerve $1$-skeleton, with weights $w_{ij}$ from overlap sizes
(dimension- and size-aware) rather than kNN adjacency. More: the chart
\emph{size} is itself a field --- larger charts average more (lower variance, more
bias) --- and it should be co-regularized with the dimension over $\nerve$, with
the same edge-preserving ($\ell_1$) penalty so that size, like dimension, stays
constant within a stratum and jumps at boundaries.

\subsection{One objective over the complex}
Attach to each chart its local data --- predicted surface, covariance, dimension,
size. The nerve records which charts overlap. The global estimate minimizes
\[
  \sum_{c}\ \underbrace{L_c(\text{local fit})}_{\text{fidelity}}
  \;+\;\sum_{\sigma\in\nerve}\ \omega_\sigma\,
  \underbrace{D_\sigma(\text{disagreement on }{\textstyle\bigcap_{c\in\sigma}U_c})}_{\text{gluing}},
\]
where $\sigma$ ranges over the simplices (edges, triangles, \dots) of $\nerve$ and
$\omega_\sigma$ is the overlap weight. PS-LPS prediction synchronization, the
covariance fusion, and the dimension/size field are all special cases --- each is
a particular choice of what the local data is and how disagreement on an overlap
is measured. (Formally this is a cochain energy of a cellular sheaf over $\nerve$;
we use that only as the name for ``make the local data agree on overlaps,'' not for
any topological content.)

\begin{plain}
Every layer of the program is, underneath, the same sentence: each chart fits its
own little patch of data, and overlapping charts are pulled toward agreeing where
they overlap. The nerve is just the precise list of overlaps --- including the
three-way ones --- and how big each is. Writing the program against that list,
instead of against a generic neighbor graph, makes every layer use the same,
geometry-aware bookkeeping.
\end{plain}

%==============================================================================
\section{What makes a complex ``good'' here: a statistical criterion}\label{sec:criterion}
%==============================================================================

If we are not chasing topology, what tells us we have the \emph{right} complex?
The answer is the only one that matters for the program: a complex is good to the
extent that estimates glued over it predict held-out data better. The overlap
structure is a modeling choice --- which local estimates to tie together, and how
strongly --- and modeling choices in this program are settled by held-out
conditional-expectation accuracy (the same rule the implementation notes use to
select the dimension field's $\lambda$). This dissolves the would-be
chicken-and-egg of ``you need the geometry to build the cover, and the cover to
estimate the geometry'': we do not need the cover to be correct in any absolute
sense, only useful, and usefulness is measurable.

It also turns the user's original observation --- that the substrate should be
\emph{derived from} the chart models --- into a concrete loop
(Figure~\ref{fig:fig_refine_loop.pdf}). Start from charts (each with a size and a
local dimension); take the nerve of their cover; glue the local estimates over it;
score the result by held-out prediction error; and let that score re-size and
re-dimension the charts. A fixed point of this loop is a self-consistent atlas
whose overlaps were chosen by the estimation problem itself. In practice one would
run a few rounds rather than seek an exact fixed point; whether the loop converges,
and to what, is an open question (Section~\ref{sec:open}), but each round is
well-defined and each is judged by accuracy.

\fig{fig_refine_loop.pdf}{0.74\linewidth}{The estimation problem chooses the
substrate. Charts (with their sizes and local dimensions) induce a cover; its
nerve says who overlaps; local estimates are glued over the nerve into
conditional-expectation and covariance fields; held-out accuracy scores the result
and feeds back to re-size and re-dimension the charts. The cover is selected, not
fixed.}

%==============================================================================
\section{The compositional case: part of the nerve is known for free}\label{sec:compositional}
%==============================================================================

For 16S relative-abundance data the nerve is not entirely something to estimate ---
much of its coarse combinatorics is handed to us by the simplex. A sample's
support (which phylotypes are present) places it on a face; two charts on
different faces $F_A$ and $F_{A'}$ can overlap only through their shared sub-face,
and which faces are adjacent is fixed by the face lattice (the support patterns
enumerated by experiment E1 of the implementation notes). So the cross-face part of
$\nerve$ --- which strata are even capable of being coupled --- is a known, finite
poset, not a fitted object.

What remains to estimate is the \emph{within-face} part: among samples sharing a
support, which charts overlap, at what sizes, and where the effective dimension
drops (the gaps measured by E2). This is exactly where the variable-size,
dimension-aware nerve earns its keep, and it scopes the work honestly: the
expensive, learned part of the substrate lives inside the strata and at their soft
boundaries (near-ties, near-zeros), while the skeleton across strata is read off
the zero pattern.

\begin{plain}
In microbiome data, the big pieces of the ``who can overlap whom'' map are free:
two samples made of completely different microbes simply do not share a chart.
What we actually have to learn is the fine structure \emph{within} a group of
samples that share the same microbes --- how their charts overlap and how many
directions they really spread in.
\end{plain}

%==============================================================================
\section{Costs, and a minimal-first path}\label{sec:plan}
%==============================================================================

\paragraph{Costs, stated plainly.} Higher simplices are combinatorially heavier:
the number of triple overlaps can be large, so the complex is truncated at
dimension two or three and built only among nearby charts --- still a sparse object,
but a real implementation cost. Variable chart size adds a selection problem (more
adaptivity means more to choose and more variance), which is why the size field
must be regularized jointly with the dimension over $\nerve$. And the assembly and
cross-validation bookkeeping must be re-expressed against the nerve rather than the
kNN graph; because PS-LPS already assembles its system from chart overlaps, this is
a generalization of existing code, not a rewrite.

\paragraph{Minimal-first.} The idea should be earned one measurable step at a time,
on the synthetic generators of the implementation notes (where the true
local-dimension field, and hence the intended overlaps, are known):
\begin{enumerate}[leftmargin=1.6em]
\item Replace PS-LPS's pairwise overlap graph with the nerve $1$-skeleton of
\emph{variable-size} charts. Test whether this alone improves held-out
Truth-RMSE on a stratified generator. This is a small change with an immediate
adaptivity payoff.
\item Add $2$-simplices and higher-order synchronization on triple overlaps. Test
whether they measurably help --- the hypothesis is that they matter only near
junctions of strata.
\item Co-regularize chart \emph{size} and chart \emph{dimension} as fields over
$\nerve$ with an $\ell_1$ penalty, so both stay constant within a stratum and jump
at boundaries.
\item Only if steps 1--3 pay off, lift prediction, covariance, and dimension into
the single cochain objective of Section~\ref{sec:gains}.
\end{enumerate}
At every step the substrate is judged by held-out estimation accuracy, never by a
topological or intrinsic-geometry score.

%==============================================================================
\section{Open questions}\label{sec:open}
%==============================================================================

\begin{itemize}[leftmargin=1.4em]
\item \textbf{Sizing for a statistically good cover.} How should chart sizes be
chosen so that the nerve is a useful ledger --- jointly with the local dimension,
and judged by accuracy? This is the substantive modeling question.
\item \textbf{The refinement loop.} Does the charts~$\to$~nerve~$\to$~glue~$\to$~refine
loop converge, and is the fixed point meaningful, or is a fixed number of rounds
the right thing?
\item \textbf{Does higher order pay?} The empirical crux: do $2$-simplex (triple)
constraints measurably improve estimates over the $1$-skeleton, and only at
junctions, as conjectured?
\item \textbf{Cheap truncation.} What is the cheapest construction of the
dimension-$\le 2$ nerve that keeps the overlaps that matter for estimation?
\item \textbf{Relation to existing machinery.} How does a nerve-based assembly sit
with the metric-graph and graph-trend-filtering tooling already in
\texttt{geosmooth} --- can the same solvers be reused on the nerve $1$-skeleton?
\end{itemize}

These are the right next questions precisely because they are all answerable by
experiment on synthetic data with known truth, and then by held-out accuracy on the
VALENCIA-shaped generators --- which keeps the whole idea, ambitious as it is,
anchored to biology and to measurable estimation gains rather than to topology.

%==============================================================================
\begin{thebibliography}{9}\small
\bibitem{impl} \emph{LPS Local Dimension Regularization: Implementation Notes},
  \nolinkurl{~/current_projects/geosmooth/project_briefs/lps_local_dimension_regularization_implementation_notes_2026-06-10}.
\bibitem{expo} \emph{Local Dimension Regularization in Local Polynomial
  Smoothing}, \nolinkurl{~/current_projects/geosmooth/project_briefs/lps_local_dimension_regularization_2026-06-10}.
\bibitem{chartaware} \emph{Chart-Aware Local Association Structure on the
  $L^\infty$ Atlas},
  \nolinkurl{~/current_projects/geosmooth/project_briefs/chart_aware_local_association_design_2026-06-09}.
\bibitem{plan} \emph{LPS Correctness and Validation: A Precise Experimental Plan},
  \nolinkurl{~/current_projects/geosmooth/project_briefs/lps_experimental_plan_2026-06-09}.
\bibitem{mapper} G.~Singh, F.~M\'emoli, G.~Carlsson, Topological methods for the
  analysis of high dimensional data sets and 3D object recognition (the Mapper
  algorithm), \emph{Eurographics SPBG} (2007).
\bibitem{sheaf} J.~Hansen, R.~Ghrist, Toward a spectral theory of cellular
  sheaves, \emph{J.\ Appl.\ Comput.\ Topology} 3 (2019).
\bibitem{nerve} K.~Borsuk, On the imbedding of systems of compacta in simplicial
  complexes (the Nerve Lemma), \emph{Fund.\ Math.} 35 (1948); J.~Leray (1945).
\end{thebibliography}

\end{document}
